\chapter{Building IR Safely}
\label{ch:pebble-builders}

A construction interface is not a verifier with fewer checks.
Its job is to make ordinary construction convenient while preventing local mistakes from corrupting storage.
Some legal programs require temporary holes: the header of a loop refers to a value that its latch has not yet constructed.
The interface must admit that state without letting execution mistake it for a complete program.

\section{Construction and verified states}
\label{sec:pebble-building-state}
\code{pir\_init} creates an empty object.
\code{pir\_block} appends a block and returns its index, or the distinct failure sentinel \code{PIR\_BLOCKS}.
\code{pir\_value} appends a value in an existing, unterminated block and returns its nonzero ID.
\code{pir\_term} sets or replaces the block terminator.
The final operation of construction is \code{pir\_verify}.

The two ID spaces deliberately have different failure conventions.
Zero is a valid block ID and an invalid value ID.
A client that writes \enquote{if the block ID is zero, fail} rejects the entry block and may accept the actual failure sentinel.
The tests check block identities explicitly rather than converting IDs to booleans.

\section{Build the running expression}
\label{sec:pebble-build-example}
Algorithm~\ref{alg:pebble-build-expression} constructs the specimen in source evaluation order.
Its seven values are input, constant 7, the first addition, constants 3 and 5, their addition, and the multiplication.
The return consumes value 7.
There is no separate copy instruction for each source subexpression because each computed value already has an ID.

\begin{algorithm}[htbp]
\caption{Building the running expression with checked allocation.}
\label{alg:pebble-build-expression}
\begin{pseudocode}
\Require fresh temporary function $F$
\Ensure verified expression function, or failure before publication
\State create entry block 0, checking its block sentinel
\State $v_1\gets input$; $v_2\gets const\ 7$
\State $v_3\gets add\ v_1,v_2$
\State $v_4\gets const\ 3$; $v_5\gets const\ 5$
\State $v_6\gets add\ v_4,v_5$; $v_7\gets mul\ v_3,v_6$
\State reject if any returned value ID is zero
\State terminate block 0 with $ret\ v_7$
\State verify $F$ before publishing it
\end{pseudocode}
\end{algorithm}

The source compiler in the repository generalizes this pattern with recursive descent.
It always emits input first, even when the expression is constant.
Dead-code elimination later removes that unused input.
The front end therefore has one simple rule instead of an early special case for every constant-only program.

\section{Backpatching a loop}
\label{sec:pebble-backpatch}
The loop fixture first allocates blocks for entry, header, body, and exit.
It creates a header phi with the entry operand filled and the latch operand temporarily zero.
It then creates the body increment, fills the pending operand, increments the revision, and verifies the complete function.
No analysis is published while that hole exists.
This is \term{backpatching}: recording an obligation now and resolving it when the referenced object becomes available.

The public teaching structs allow the fixture to perform the backpatch directly.
A production library should hide that access behind a construction token or an edit session that records the pending edges and refuses to seal while obligations remain.
Do not claim that the current builders make every representable C object valid; they make bounded construction possible, and the verifier establishes validity afterward.

\section{Diagnostics as an interface}
\label{sec:pebble-diagnostics}
Failures return a status and a bounded diagnostic string.
The current messages identify a category such as \enquote{operand does not dominate its use}.
They do not yet carry source spans, value IDs, or an explanation path through the graph.
That is sufficient for the tests, but a user-facing compiler should attach the failed value and block, retain source locations from construction, and display the original use alongside its unavailable definition.

Local allocation is constant time.
A successful expression build takes linear time in syntax-tree size before verification, with recursion bounded by the parser's explicit nesting limit.
Whole-function transaction copying costs the function object's storage size, including unused capacity in this fixed-array version.
That copying cost is a deliberate price for failure atomicity.

\section{Exercises}
\label{sec:pebble-six-exercises}
\begin{enumerate}
\item \label{ex:pebble-builder-order} Why does appending a value after installing a terminator fail, although replacing that terminator is permitted?
\item \label{ex:pebble-pending-phi} A loop is verified before its latch operand has been backpatched. Which invariant rejects it, and why must a zero operand not be interpreted as constant zero?
\end{enumerate}

\section{Further reading}
\label{sec:pebble-six-reading}
Appel's treatment of intermediate-code construction gives useful context for separating translation bookkeeping from later analysis~\cite{appel-1998-modern}.
